% Generated draft: requires mathematical and source review.
\documentclass[11pt]{article}
\usepackage[margin=25mm]{geometry}
\usepackage{fontspec}
\setmainfont{Latin Modern Roman}
\usepackage{amsmath,amssymb,mathtools,unicode-math}
\setmathfont{Latin Modern Math}
\usepackage{enumitem,needspace,xurl,hyperref}
\setlength{\parindent}{0pt}
\setlength{\parskip}{5pt}
\setlength{\emergencystretch}{3em}
\newcommand{\N}{\mathbb{N}}
\newcommand{\Z}{\mathbb{Z}}
\newcommand{\Q}{\mathbb{Q}}
\newcommand{\R}{\mathbb{R}}
\newcommand{\C}{\mathbb{C}}
\newcommand{\F}{\mathbb{F}}
\newcommand{\A}{\mathbb{A}}
\newcommand{\PP}{\mathbb{P}}
\newcommand{\E}{\mathbb{E}}
\newcommand{\eps}{\varepsilon}
\newcommand{\dd}{\,\mathrm{d}}
\begin{document}
\section*{UnsolvedMath ID: 30001604}
\textbf{Source problem code:} catalogue entry
\subsection{Definitions and mathematical statement}
Let \(S\) denote the smooth surface in the supplied statement. The source does not specify its base field, characteristic, properness, or projectivity. The following definitions give a conventional framework for interpreting the question; the source does not uniquely prescribe this framework.

After fixing a field \(k\) and interpreting \(S\) as a smooth surface over \(k\), a finite zero-dimensional subscheme \(Z\subset S\) has length \(\ell(Z)=\dim_k H^0(Z,\mathcal O_Z)\). A conventional one-parameter collision is described by a discrete valuation \(k\)-algebra \(A\), with fraction field \(K\) and residue field \(\kappa\), and a finite flat closed subscheme
\[
\mathcal Z\subset S\times_k\operatorname{Spec}A.
\]
Its generic fibre \(\mathcal Z_K\) consists of finite subschemes with pairwise disjoint supports, while some of these supports coalesce in the special fibre \(\mathcal Z_\kappa\). The latter is called a collision limit. Finite flatness preserves the total fibre length. Neither the allowed generic subschemes nor the permitted collision families are specified in the supplied record. A fat point of multiplicity \(m\geq1\) at a closed point \(p\) is the subscheme defined by \(\mathcal I_p^m\), where \(\mathcal I_p\) is the ideal sheaf of \(p\). The accompanying literature assessment discusses fat points, whereas the statement itself concerns zero-dimensional schemes generally.

For a line bundle \(L\) on \(S\) and a finite-dimensional space \(V\subseteq H^0(S,L)\), define
\[
V(-Z)=\ker\bigl(V\longrightarrow H^0(Z,L|_Z)\bigr).
\]
The associated linear series of sections vanishing on \(Z\) is \(\mathbb P(V(-Z))\); when \(V(-Z)\neq0\), its dimension is \(\dim_k V(-Z)-1\). When \(V(-Z)=0\), the series is empty. This notation makes precise one possible meaning of estimating dimensions of linear series subject to zero-dimensional conditions; the source does not specify \(L\), \(V\), or the desired bounds.

\textbf{Requested construction and classification.} Construct collision limits of zero-dimensional schemes on \(S\) that permit useful estimates for dimensions of linear series subject to the original zero-dimensional conditions, and classify collisions involving a small number of points.

This is a research programme rather than a uniquely quantified assertion in the supplied record. In particular, the number of points, their lengths or multiplicities, the meaning of usefulness, the required estimates, and the equivalence relation for classification remain unspecified. It is also unclear whether every support must merge into a single point. No restriction to fat points, general points, an algebraically closed field, or a particular characteristic is imposed here as an unstated source hypothesis.
\subsection{Short English statement}
The problem asks for controlled ways of merging finite subschemes on a smooth surface so that the resulting limits help estimate dimensions of linear series, together with a classification for collisions involving few points. The supplied record leaves the allowed schemes, number of points, classification equivalence, and required estimates unspecified, so it does not determine a single precise quantified conjecture.
\subsection{Sources}
\begin{itemize}
\item Local UnsolvedMath registry: supplied record with permanent ID \(30001604\), title \emph{Collisions of Zero-Dimensional Schemes on Surfaces}, exact source problem code \texttt{catalogue entry}, in the Oberwolfach Reports Open Problems collection. Project registry locations are \texttt{data/unsolvedmath/problems.json} and the supplementary \texttt{catalog/problems.json}; their contents were not independently inspected for this draft. No individual proposer is supplied.
\item Supplied report citation: \emph{Mini-Workshop: Linear Series on Algebraic Varieties (2011)}, Oberwolfach Reports, pp.\ 2613--2650, DOI \texttt{10.4171/owr/2010/45}: \url{https://doi.org/10.4171/owr/2010/45}. The record describes Oberwolfach Reports as documenting mathematical workshops and their open-problem discussions. It separately gives proposed year \(2010\); the apparent year discrepancy is unresolved.
\item Additional supplied literature source, without a supplied title or author attribution: \url{https://doi.org/10.5802/aif.2319}.
\item Additional supplied literature source, without a supplied title or author attribution: \url{https://ems.press/content/serial-article-files/46305?nt=1}.
\item Literature assessment dated exactly \texttt{2026-08-21}: the supplied curation labels the problem open and states that the general classification of collisions of fat points remains beyond current results and that the literature identifies it as harder than the open Nagata problem. These are supplied dated curation claims, not independently verified mathematical facts or a new literature assessment. The record explicitly requests independent verification of time-sensitive status and cited claims.
\item The supplied statement carries the label \texttt{corrected\_verified} and reports a manual semantic, typographic, and LaTeX review of its cleanup. That supplied review label does not resolve the formulation gaps or constitute mathematical verification of this draft.
\end{itemize}
\end{document}
